\documentclass[10pt,a4paper]{article}
\usepackage[utf8]{inputenc}
\usepackage[a4paper,left=3cm,right=2cm]{geometry}
\usepackage{amsmath}
\usepackage[thmmarks, amsmath, thref]{ntheorem}
\usepackage{amsfonts}
\usepackage{amssymb}
\usepackage{fancyhdr}
\pagestyle{fancy}
\usepackage{anysize}
\usepackage{graphicx}
\usepackage{float}
\usepackage{hyperref}
\usepackage{enumitem}
\usepackage{verbatim}
\newtheorem{theorem}{Theorem}
\theoremsymbol{\ensuremath{\blacksquare}}
\newtheorem*{solution}{Solución:}

\usepackage{xcolor}

% Margenes y formalidades

\textwidth 6.25in % Margen Derecho y ancho del texto
\textheight 8.5in
\oddsidemargin 0in % Margen Izquierdo
\headheight 0in % Largo del texto
\leftmargin 1in
\parindent 0em % Salto entre la indentación (Sangría)https://www.overleaf.com/4519127669sdpptkhvsbfn
\topmargin -0.5in % Margen Superior
\parskip 0.5ex % Salto entre parrafos
\baselineskip 1pt



\lhead{
\includegraphics[scale=0.08]{Figuras/logoUSM.png}} % texto izquierda de la cabecera
%\chead{}                                      % texto centro de la cabecera
\rhead{Universidad Técnica Federico Santa María\\Ingeniería Comercial } % número de página a la derecha
\renewcommand{\headrulewidth}{0.4pt} % grosor de la línea de la cabecera
\renewcommand{\footrulewidth}{0.4pt} % grosor de la línea del pie

\begin{document}
\vspace{5cm}



\begin{center}
\begin{large}
\textbf{\\CERTAMEN 1:}\textbf{ ECONOMETRÍA ICS-294}
\end{large}\\
\vspace{0.2cm}
II Semestre 2025\\
Profesores: Jocelyn Tapia Stefanoni, Franco Antonucci Fernández y Ricardo Pruneda Urz\'ua\\
%Ayudante: 
\end{center}

\textbf{Instrucciones:} 

\begin{itemize}

    \item Usted dispone de 90 minutos  para intentar los 100 puntos del certamen.
    \item Cuide la presentación y redacción de sus preguntas.
    \item No puede utilizar celular, computador o tablet, pero sí una calculadora TI.
     
\end{itemize}

\section{Comentes y preguntas cortas (40 puntos)}

Responda solo 4 de las siguientes preguntas: 
\begin{enumerate}
   \item \textbf{RdeA9.1} (10 puntos) 
¿Cuál es la diferencia entre un muestreo aleatorio simple y un muestreo estratificado? ¿En qué situaciones conviene usar cada uno en una investigación econométrica?
 
{\color{blue}(2,5 puntos) Muestreo aleatorio simple 

Cada individuo de la población tiene la misma probabilidad de ser seleccionado. 
Se seleccionan $n$ observaciones de manera completamente aleatoria, sin importar sus características.

(2,5 puntos) Muestreo estratificado

La población se divide primero en estratos homogéneos (grupos según alguna característica relevante: 
edad, género, nivel educativo, región, sector económico, etc.). 
Luego se selecciona una muestra aleatoria dentro de cada estrato.

¿Cuándo conviene usar cada uno?

(2,5 puntos) Muestreo aleatorio simple
\begin{itemize}
    \item Cuando la población es relativamente homogénea respecto a las variables de interés.
    \item Cuando no se dispone de información previa para clasificar en estratos.
  
\end{itemize}

(2,5 puntos) Muestreo estratificado
\begin{itemize}
    \item Cuando la población es heterogénea y existen subgrupos claramente diferenciados.
    \item Si se quiere garantizar representatividad de cada estrato (por ejemplo, hombres y mujeres, sectores productivos).
    \item Reduce la varianza de los estimadores y mejora la precisión de los resultados.
\end{itemize}}

    \item 
En una importante empresa exportadora, el gerente de logística sostiene que el tiempo promedio de entrega de los pedidos es de 15 días.  
Para verificarlo, una analista toma una muestra aleatoria de 400 envíos y encuentra que el tiempo promedio de entrega es de 14,26 días, con una desviación estándar de 2 días.  
Con un 95\% de confianza, ¿se puede concluir que la afirmación del gerente es correcta?


{\color{blue}
Queremos contrastar si el tiempo promedio de entrega poblacional es $\mu_0=15$ d\'ias.
A partir de una muestra aleatoria de tama\~no $n=400$, se obtiene media muestral $\bar{x}=14.6$ y desviaci\'on est\'andar muestral $s=2$.


Como $n$ es grande, usamos aproximaci\'on normal con error est\'andar
\[
SE=\frac{s}{\sqrt{n}}=\frac{2}{\sqrt{400}}=\frac{2}{20}=0.1.
\]
(5 puntos)

\textbf{ Intervalo de confianza (95\%) para $\mu$}
\[
\bar{x}\pm z_{0.975}\,SE
=14.26\pm 1.96(0.1)
=14.26\pm 0.196
=(14.064,\ 14.456).
\]
El valor $15$ no pertenece a este intervalo. 
Con un 95\% de confianza, la evidencia emp\'irica \textbf{refuta} la afirmaci\'on del gerente: el tiempo promedio de entrega es
\emph{significativamente menor} que 15 d\'ias. (5 puntos)
}

    \item \textbf{RdeA3.2} (10 puntos) 
    Un reciente estudio propone que el {\bf uso de redes} sociales influyen en el rendimiento de los estudiantes. Se estima el siguiente  modelo:
    
    $$y_i = \hat{\beta}_0 + \hat{\beta}_1 \cdot x_i + \hat{u}_i$$
    Donde $y_i$ es el promedio de notas del alumno $i$; $x_i$ es la cantidad de aplicaciones instaladas en el celular del alumno $i$. $\hat{\beta}_0$; $\hat{\beta}_1$; $u$ son los coeficientes estimados.\\
    
    En ausencia de variable omitida, este modelo entregará un coeficiente insesgado. Comente.\\

    {\color{blue}\textbf{Respuesta:} Falso o incierto. Se podría sospechar que la variable independiente tiene un error de medición, ya que lo que importa no es la cantidad de aplicaciones instaladas sino cuántas utiliza \textbf{(5 puntos)}, por lo tanto como el error está en la variable independiente, se tiene:
    $$y_i = b_0 + b_1 \cdot x^{'}_i + v$$

    Donde ahora $x^{'} = x + u$. Luego, $v = e + \beta_1 \cdot u$. Desarrollando se llega a que \textbf{(2,5 puntos)}
    $$E[\hat{b}_1] = \beta_1 \cdot \frac{var(x)}{var(x) + var(u)}$$
    
    Por lo tanto, el coeficiente tiene sesgo \textbf{(2,5 puntos)}.
}

    \item (10 puntos)
Considere el modelo de regresión lineal simple:

\[
y_i = \beta_0 + \beta_1 x_i + u_i, \qquad i=1,\dots,n
\]

donde $E[u_i|x_i]=0$ y $Var(u_i|x_i)=\sigma^2$.

Muestre que el estimador de Mínimos Cuadrados Ordinarios para la pendiente:
    \[
    \hat{\beta}_1 = \frac{\sum_{i=1}^n (x_i - \bar{x})(y_i - \bar{y})}{\sum_{i=1}^n (x_i - \bar{x})^2}.
    \]

Es insesgado y su varianza es:
    \[
    Var(\hat{\beta}_1) = \frac{\sigma^2}{\sum_{i=1}^n (x_i - \bar{x})^2}.
    \]



El estimador de MCO para $\beta_1$ es:
\[
\hat{\beta}_1 = \frac{\sum_{i=1}^n (x_i - \bar{x})(y_i - \bar{y})}{\sum_{i=1}^n (x_i - \bar{x})^2}.
\]

{\color{blue} Sustituyendo $y_i = \beta_0 + \beta_1 x_i + u_i$:
\[
\hat{\beta}_1 = \frac{\sum_{i=1}^n (x_i - \bar{x})(\beta_0 + \beta_1 x_i + u_i - \bar{y})}{\sum_{i=1}^n (x_i - \bar{x})^2}.
\]
(1 puntos)

Se puede demostrar que los términos constantes se cancelan y obtenemos:
\[
\hat{\beta}_1 = \beta_1 + \frac{\sum_{i=1}^n (x_i - \bar{x})u_i}{\sum_{i=1}^n (x_i - \bar{x})^2}.
\]
(2 puntos)

Tomamos esperanza condicional en $x$:
\[
E[\hat{\beta}_1 \mid x] = \beta_1 + 
\frac{\sum_{i=1}^n (x_i - \bar{x})E[u_i \mid x_i]}{\sum_{i=1}^n (x_i - \bar{x})^2}.
\]
 
Por el supuesto de media condicional cero, $E[u_i|x_i]=0$, entonces:
\[
E[\hat{\beta}_1 \mid x] = \beta_1.
\]
 (2 puntos)
Por lo tanto:
\[
E[\hat{\beta}_1] = \beta_1,
\]
es decir, $\hat{\beta}_1$ es un estimador \textbf{insesgado}.



De la expresión anterior:
\[
\hat{\beta}_1 - \beta_1 = \frac{\sum_{i=1}^n (x_i - \bar{x})u_i}{\sum_{i=1}^n (x_i - \bar{x})^2}.
\]

Por lo tanto:
\[
Var(\hat{\beta}_1 \mid x) = 
\frac{Var\left(\sum_{i=1}^n (x_i - \bar{x})u_i \;\middle|\; x \right)}{\left(\sum_{i=1}^n (x_i - \bar{x})^2\right)^2}.
\]

Dado que los $u_i$ son homocedásticos e independientes con $Var(u_i|x_i)=\sigma^2$: (2 puntos)
\[
Var\left(\sum_{i=1}^n (x_i - \bar{x})u_i \;\middle|\; x \right) 
= \sum_{i=1}^n (x_i - \bar{x})^2 \sigma^2.
\]

Sustituyendo:
\[
Var(\hat{\beta}_1 \mid x) = \frac{\sigma^2 \sum_{i=1}^n (x_i - \bar{x})^2}{\left(\sum_{i=1}^n (x_i - \bar{x})^2\right)^2}.
\]
(2 puntos)
Simplificando:
\[
Var(\hat{\beta}_1) = \frac{\sigma^2}{\sum_{i=1}^n (x_i - \bar{x})^2}.
\] (1 punto)
}



\item \textbf{RdeA9.1} (10 puntos) 
Considere un modelo de regresión lineal en el que se cumple la descomposición:

\[
Y_i = \hat{Y}_i + \hat{u}_i, \qquad i = 1,\dots,n,
\]

donde $Y_i$ es el valor observado de la variable dependiente, $\hat{Y}_i$ es el valor ajustado por el modelo y $\hat{u}_i$ es el residuo. 

Se propone la siguiente igualdad:
\[
\sum_{i=1}^n (Y_i - \bar{Y})^2 
= \sum_{i=1}^n (\hat{Y}_i - \bar{Y})^2 
+ \sum_{i=1}^n \hat{u}_i^2.
\]


\begin{enumerate}
    \item 
    \text{(4 puntos) Demuestre o refute esta igualdad mostrando los pasos necesarios.}  

Partimos de la descomposición clásica del modelo de regresión:
\[
Y_i = \hat{Y}_i + \hat{u}_i.
\]

Restando la media \(\bar{Y}\) a ambos lados, obtenemos:
\[
Y_i - \bar{Y} = (\hat{Y}_i - \bar{Y}) + \hat{u}_i.
\]

Elevando al cuadrado y sumando sobre \(i=1,\ldots,n\):
\[
\sum_{i=1}^n (Y_i - \bar{Y})^2 = \sum_{i=1}^n \left[ (\hat{Y}_i - \bar{Y}) + \hat{u}_i \right]^2.
\]

Expandiendo el cuadrado:
\[
\sum_{i=1}^n (Y_i - \bar{Y})^2 = \sum_{i=1}^n (\hat{Y}_i - \bar{Y})^2 + \sum_{i=1}^n \hat{u}_i^2 + 2 \sum_{i=1}^n (\hat{Y}_i - \bar{Y}) \hat{u}_i.
\]

Por lo tanto, la igualdad propuesta se cumple si y sólo si:
\[
 \sum_{i=1}^n (\hat{Y}_i - \bar{Y}) \hat{u}_i = 0,
\]

 \item   \text{(3 puntos) ¿Qué propiedad o condición debe cumplirse para que esta igualdad sea válida?} 

\[
\sum_{i=1}^n (\hat{Y}_i - \bar{Y}) \hat{u}_i
\]

El primer término es \(0\) por la ortogonalidad de X y $\hat{u}$, como \(\hat{Y}=X\hat{\beta_1}\), entonces la suma será 0 (por las condiciones de primer orden de MCO).

El segundo término es \(0\) porque la suma de los residuales es cero.


 \item   \text{(3 puntos) Explique el significado de cada término en esta expresión.} 

- \(\displaystyle \sum_{i=1}^n (Y_i - \bar{Y})^2\): \textit{Suma total de cuadrados} (SST), representa la variabilidad total de la variable dependiente \(Y\) respecto a su media.

- \(\displaystyle \sum_{i=1}^n (\hat{Y}_i - \bar{Y})^2\): \textit{Suma de cuadrados explicada} (SSR), mide la variabilidad de los valores predichos por el modelo respecto a la media, es decir, la parte de la variabilidad que el modelo logra explicar.

- \(\displaystyle \sum_{i=1}^n {u}_i^2\): \textit{Suma de cuadrados de residuos} (SSE), representa la variabilidad no explicada por el modelo, es decir, la variabilidad de los  residuos.
\end{enumerate}

    \item \textbf{RdeA3.2} (10 puntos) Suponga que se tiene el siguiente modelo 
    $$\text{salario} = \beta_0 + \beta_1 \cdot \text{educación}  + u$$
    y, al mismo tiempo se tiene que 
    $$\text{educación} = \gamma_0 + \gamma_1 \cdot \text{salario} + v$$
    
    Como la relación entre salario y educación es positiva, entonces se tendrá que $\beta_1 = \gamma1$. Comente.\\

\textbf{Respuesta:} Falso o incierto. Se puede intuir de buena manera que las variables correlacionan estadísticamente de forma positiva, pero la causalidad económica nos dice que mientras más años de educación mayor salario, pero no necesariamente un mayor salario causa mayor cantidad de años de educación \textbf{(5 puntos)}. Es más, matemáticamente se tiene que

$$\beta_1 = \frac{cov(\text{educación}, \text{salario})}{var(\text{educación})}$$
$$\gamma_1 = \frac{cov(\text{salario}, \text{educación})}{var(\text{salario})}$$

Luego $\beta_1 \neq \gamma_1$. \textbf{(5 puntos)}



\end{enumerate}

\section{Ejercicios (60 puntos)}

Responda solo 2 de las siguientes preguntas.
\begin{enumerate}
      \item \textbf{RdeA3.2} (30 puntos)
      Suponga que usted quiere estudiar el impacto que tiene la escolaridad de los padres en la escolaridad de los hijos, y estima el siguiente modelo:

\begin{equation}
    S_i = \alpha + \beta S^P_i + \mu_i
\end{equation}


      Donde $S_i$ es la escolaridad del hijo $i$; $S^P_i$ representan los años de escolaridad promedio de los padres del individuo $i$; mientras que $\mu_i$ es un término de error bien comportado. Lamentablemente usted no incluye la habilidad de los hijos como determinante clave de su escolaridad. En realidad, la verdadera relación existente entre escolaridad de los hijos y la de sus padres debería ser estimada a través de la siguiente ecuación:
\begin{equation}
    S_i = \alpha + \beta S^P_i + \gamma H_i + \mu_i
\end{equation}

        \begin{itemize}
            \item [a)] (10 puntos) Demuestre que al estimar $\beta$ a partir del modelo univariado se incurre en sesgo.\\

            \textbf{Respuesta:} El estimador de MCO de $\beta$, para el modelo univariado viene dado por:

            $$\hat{\beta} = \frac{cov(S^P_i, S_i)}{var(S^P_i)}$$

            Por otro lado, si se reemplaza el valor de $S_i$ de la ecuación (2) se obtiene:

            $$\hat{\beta} = \frac{cov(S^P_i, \beta S^P_i + \gamma H_i + \mu_i)}{var(S^P_i)}$$

            Es decir,

            $$\hat{\beta} = \beta + \gamma\frac{cov(S^P_i, H_i)}{var(S^P_i)} + \frac{cov(S^P_i,\mu_i)}{var(S^P_i)}$$

            Aplicando valor esperado a ambos lados:

            $$E[\hat{\beta}] = \beta + \gamma\frac{cov(S^P_i, H_i)}{var(S^P_i)}$$

            Por lo tanto, estimar $\beta$ a través del modelo univariado provoca un sesgo, el cuasl viene dado por
            $$E[\hat{\beta}] - \beta = \gamma\frac{cov(S^P_i, H_i)}{var(S^P_i)}$$


            \item [b)] (5 puntos) ¿De qué depende el sesgo de la estimación?\\

            \textbf{Respuesta:} Como puede apreciarse en el resultado anterior, el sesgo depende de dos elementos:
            \begin{itemize}
                \item $\gamma$: El efecto que tiene la variable omitida, habilidad, sobre la variable dependiente del modelo (escolaridad)
                \item $cov(S^P_i, H_i)$ la relación que existe entre la variable omitida, habilidad, y la variable explicativa incluida en la estimación del modelo (escolaridad de los padres).
            \end{itemize}
 
            \item [c)] (5 puntos) Es posible especular respecto de la dirección del sesgo?\\

            \textbf{Respuesta:} Sí. En primer lugar es muy posible que el efecto que tiene la habilidad sobre la escolaridad de las personas sea positivo ($\gamma$), ya que personas más hábiles debieran alcanzar, en promedio, mayores niveles educativos. Por otro lado, es también más probable que padres con mayor escolaridad tengan a su vez hijos con mayor habilidad (hay un mayor estímulo y, además, es probable que exista una correlación positiva entre habilidades intergeneracionales). Por lo tanto, es muy probable que el estimador esté sesgado hacia arriba.

            \item [d)] (10 puntos) Si usted cuenta con una muestra de 5 individuos:

            \begin{table}[H]
                \centering
                \begin{tabular}{c|c|c}
                 Escolaridad ($S$) & Escolaridad de los padres ($S^P$)& Habilidad ($H$) \\ \hline
                 8&4&6\\  
                 10&13&5\\ 
                 12&13&8\\ 
                 17&15&10\\ 
                 22&21&11\\ 
                \end{tabular}
                %\caption{Caption}
                \label{tab:placeholder}
            \end{table}

            Sin considerar habilidad, ¿cuál sería el coeficiente estimado para $\beta$? Y si se incluye habilidad, cómo debería ser el coeficiente estimado para $\beta$?\\
            
            \textbf{Respuesta:} De los datos se obtiene que:

            $$\hat{\beta} = \frac{cov(S^P_i, S_i)}{var(S^P_i)}=\frac{31,05}{37,2} \approx 0,84$$

            Además, si se incluye la habilidad lo más probable es que el estimador esté sesgado hacia arriba por lo que se esperaría que la nueva medición arrojara un coeficiente menor al encontrado anteriormente. De hecho, este se puede justificar con que habilidad con escolaridad correlacionan positivo ($\gamma >0$). Por otro lado, se esperaría que la relación entre habilidad y escolaridad de los padres fuese positiva, lo que de hecho se calcula como
            $$cov(S^P_i, H_i) \approx 11,5$$
            Por lo tanto, el sesgo de la estimación sería positivo.
        \end{itemize}

        \item (30 puntos) La dirección de un banco está interesada en conocer el impacto de un programa de mentorías internas sobre el desempeño salarial de sus empleados. 
Se ofreció la posibilidad de participar de manera voluntaria a todos los trabajadores con al menos 3 años de experiencia. 
De un total de 200 empleados, 80 participaron en el programa. 
Un año después, el salario promedio de los empleados que recibieron la mentoría es de US\$2.400, mientras que el de aquellos que no participaron es de US\$2.100.

\begin{enumerate}
    \item (8 puntos) ¿Podemos concluir que existe una relación \textit{causal} entre participar en el programa de mentorías y el aumento del salario? 
    \\[0.25cm]
   {\color{blue} No necesariamente. La participación fue \textbf{voluntaria}, por lo que es posible que los trabajadores que se inscribieron en la mentoría ya tuvieran características distintas a los que no lo hicieron 
    (por ejemplo, mayor motivación, productividad o ambición). 
    Estas características no observadas pueden estar capturadas en el término de error, generando \textbf{sesgo de selección}. 
    Por lo tanto, la diferencia observada en salarios refleja tanto el posible efecto del programa como las diferencias preexistentes entre los grupos.}

    \item (8 puntos) Calcule un intervalo de confianza al 95\% para la diferencia de medias salariales entre empleados que participaron en la mentoría y los que no participaron, 
suponiendo que la varianza poblacional de los salarios es conocida e igual a $\sigma^2 = 10000$ en ambos grupos. 
El número de empleados en el grupo de mentoría es $n_1 = 80$ y en el grupo de control es $n_0 = 120$.

{\color{blue}
La diferencia de medias es:
\[
\hat{\Delta} = \overline{Y}_{mentoría=1} - \overline{Y}_{mentoría=0} = 2400 - 2100 = 300
\] (2 puntos)

El error estándar de la diferencia es:
\[
SE = \sqrt{\frac{\sigma^2}{n_1} + \frac{\sigma^2}{n_0}} 
= \sqrt{\frac{10000}{80} + \frac{10000}{120}} 
= \sqrt{125 + 83.3} \approx 14.43
\]
(2 puntos)

Como $\sigma^2$ es conocida, usamos la distribución normal estándar $Z_{\alpha/2}$ para un nivel de confianza del 95\%:
\[
Z_{0.025} = 1.96
\]

El intervalo de confianza es:
\[
IC_{95\%} = \hat{\Delta} \pm Z_{0.025} \cdot SE 
= 300 \pm 1.96 \cdot 14.43 
= 300 \pm 28.3
\]

\[
IC_{95\%} = (271.7 \;,\; 328.3)
\]

Con un 95\% de confianza, la verdadera diferencia en los salarios promedio entre empleados que participaron en la mentoría y los que no lo hicieron 
se encuentra entre US\$272 y US\$328. (4 puntos)  }
 \item (6 puntos) Halle el valor del coeficiente $\hat{\beta}_1$ en un modelo de regresión lineal simple que explique el salario en función de la variable dicotómica 
    ``mentoria'' (1 = participó, 0 = no participó). 

{\color{blue}  El modelo estimado es:
    \[
    salario_i = \hat{\beta}_0 + \hat{\beta}_1 \cdot mentoria_i + u_i
    \]

    El coeficiente $\hat{\beta}_1$ corresponde a la diferencia de medias entre los grupos:
    \[
    \hat{\beta}_1 = \overline{Y}_{mentoria=1} - \overline{Y}_{mentoria=0} = 2400 - 2100 = 300
    \]

 Los empleados que participaron en el programa de mentorías tienen en promedio un salario 
    US\$300 mayor que los que no participaron.}
\item (8 puntos) El coeficiente $\hat{\beta}_1$ estimado, ¿representa el ATT o el ATE? Justifique.
    
    {\color{blue}En este caso, $\hat{\beta}_1$ es simplemente la \textbf{diferencia de medias observada} entre los participantes y los no participantes. 
    Como la asignación al tratamiento fue voluntaria y no aleatoria, no podemos interpretar directamente este coeficiente como un ATE ni como un ATT en sentido estricto.   (6 puntos)
    Sin embargo:
    \begin{itemize}
        \item Si hubiera una asignación aleatoria, $\hat{\beta}_1$ sería tanto el ATE como el ATT.
        \item Con asignación voluntaria, lo más cercano es que $\hat{\beta}_1$ intente capturar el ATT (efecto promedio sobre los tratados), pero contaminado por \textbf{sesgo de selección}.
    \end{itemize}
(2 puntos)    
    }
\end{enumerate}
%     \item (30 puntos) regresion calculo estimadores residuos grafico modelo de regresion  

% Considere el siguiente modelo de regresión lineal con los supuestos estándares:

%     \[
%     Y_i = \beta_0 + \beta_1 X_{1i} + \beta_2 X_{2i} + u_i
%     \]
    
% Se dispone de la siguiente información:

% \begin{table}[H]
% \centering
% \begin{tabular}{|c|c|c|c|}
% \hline
% \textbf{Lote} & \textbf{Tiempo de producción (\(Y\))} & \textbf{Nº de operarios (\(X_1\))} & \textbf{Tiempo preparación (hrs) (\(X_2\))} \\
% \hline
% 1 & 10 & 1 & 8 \\
% 2 & 13 & 2 & 7 \\
% 3 & 17 & 3 & 3 \\
% 4 & 20 & 4 & 1 \\
% \hline
% \end{tabular}
% \caption{Datos sobre producción y uso de recursos}
% \end{table}


% Una empresa desea estimar el tiempo de producción de un lote en función del número de operarios asignados al proceso y del tiempo que toma preparar la maquinaria. Se propone el siguiente modelo:

% \[
% Y_i = \beta_0 + \beta_1 \cdot \text{Operarios}_i + \beta_2 \cdot \text{Preparación}_i + \varepsilon_i
% \]


%     \text{(5 puntos) Escriba el modelo en forma matricial.} 

% El modelo en forma matricial se escribe como:
% \[
% Y = X\beta + u,
% \]

% donde

% \[
% Y=
% \begin{bmatrix}
% 10\\
% 13\\
% 17\\
% 20
% \end{bmatrix},
% \quad
% X=
% \begin{bmatrix}
% 1 & 1 & 8\\
% 1 & 2 & 7\\
% 1 & 3 & 3\\
% 1 & 4 & 1
% \end{bmatrix},
% \quad
% \beta=
% \begin{bmatrix}
% \beta_0\\
% \beta_1\\
% \beta_2
% \end{bmatrix},
% \quad
% u=
% \begin{bmatrix}
% u_1\\
% u_2\\
% u_3\\
% u_4
% \end{bmatrix}.
% \]

% Estructura matricial:
% \[
% \begin{bmatrix}
% 10\\
% 13\\
% 17\\
% 20
% \end{bmatrix}
% =
% \begin{bmatrix}
% 1 & 1 & 8\\
% 1 & 2 & 7\\
% 1 & 3 & 3\\
% 1 & 4 & 1
% \end{bmatrix}
% \begin{bmatrix}
% \beta_0\\
% \beta_1\\
% \beta_2
% \end{bmatrix}
% +
% \begin{bmatrix}
% u_1\\
% u_2\\
% u_3\\
% u_4
% \end{bmatrix}.
% \]


%     \text{(10 puntos) Calcule los estimadores \(\hat{\beta} = (X'X)^{-1}X'Y\)}.
    
%     \text{(Realice el desarrollo completo del procedimiento utilizado para obtener los resultados) }

% \[
% X'X \;=\;
% \begin{bmatrix}
% 4 & 10 & 19\\
% 10 & 30 & 35\\
% 19 & 35 & 123
% \end{bmatrix}.
% \]

% \[
% (X'X)^{-1} \;=\;
% \begin{bmatrix}
% \frac{493}{6} & -\frac{113}{6} & -\frac{22}{3}\\[6pt]
% -\frac{113}{6} & \frac{131}{30} & \frac{5}{3}\\[6pt]
% -\frac{22}{3} & \frac{5}{3} & \frac{2}{3}
% \end{bmatrix}.
% \]

% \[
% X'Y \;=\;
% \begin{bmatrix}
% 60\\
% 167\\
% 242
% \end{bmatrix}.
% \]

% \[
% \hat\beta \;=\; (X'X)^{-1}X'Y \;=\;
% \begin{bmatrix}
% \frac{61}{6}\\[6pt]
% \frac{77}{30}\\[6pt]
% -\frac{1}{3}
% \end{bmatrix}
% =
% \begin{bmatrix}
% 10.1667\\[2pt]
% 2.5667\\[2pt]
% -0.3333
% \end{bmatrix}.
% \]


    
%     \text{(7 puntos) Interprete los coeficientes estimados en el contexto operativo.} 

% \[\hat\beta_0 \approx 10.17
% \]
% (intercepto) Representa el tiempo base de producción de un lote cuando el número de operarios y el tiempo de preparación son cero. 

% \[\hat\beta_1 \approx 2.57
% \]
% (efecto de los operarios) Indica que, manteniendo fijo el tiempo de preparación, cada operario adicional aumenta en promedio unas 2.6 horas el tiempo de producción. 

% \[\hat\beta_2 \approx -0.33
% \]
% (efecto de la preparación) Significa que, en promedio, por cada hora extra de preparación, el tiempo de producción disminuye en alrededor de 0.33 horas.


%     \text{(8 puntos) Obtenga el vector de predicción Y.} 
% \

% \[
% X=
% \begin{bmatrix}
% 1 & 1 & 8\\
% 1 & 2 & 7\\
% 1 & 3 & 3\\
% 1 & 4 & 1
% \end{bmatrix},
% \qquad
% \hat\beta=
% \begin{bmatrix}
% \frac{61}{6}\\[4pt]
% \frac{77}{30}\\[4pt]
% -\frac{1}{3}
% \end{bmatrix}.
% \]

% \[
% \hat Y \;=\; X\hat\beta \;=\;
% \begin{bmatrix}
% 1 & 1 & 8\\
% 1 & 2 & 7\\
% 1 & 3 & 3\\
% 1 & 4 & 1
% \end{bmatrix}
% \begin{bmatrix}
% \frac{61}{6}\\[4pt]
% \frac{77}{30}\\[4pt]
% -\frac{1}{3}
% \end{bmatrix}
% =
% \begin{bmatrix}
% \frac{61}{6} + \frac{77}{30} - \frac{8}{3}\\[6pt]
% \frac{61}{6} + 2\!\left(\frac{77}{30}\right) - \frac{7}{3}\\[6pt]
% \frac{61}{6} + 3\!\left(\frac{77}{30}\right) - 1\\[6pt]
% \frac{61}{6} + 4\!\left(\frac{77}{30}\right) - \frac{1}{3}
% \end{bmatrix}.
% \]

% \text{Simplificando cada componente:}
% \[
% \hat Y =
% \begin{bmatrix}
% \frac{151}{15}\\[6pt]
% \frac{389}{30}\\[6pt]
% \frac{253}{15}\\[6pt]
% \frac{201}{10}
% \end{bmatrix}
% \approx
% \begin{bmatrix}
% 10.0667\\[2pt]
% 12.9667\\[2pt]
% 16.8667\\[2pt]
% 20.1000
% \end{bmatrix}.
% \]


\item (30 puntos) Un grupo de investigadores estudia la relación entre el coeficiente intelectual (CI) y el puntaje en una prueba estandarizada. Se cuenta con 7 observaciones de individuos, para los cuales se registró la siguiente información:

\begin{center}
\begin{tabular}{|c|c|c|}
\hline
\textbf{Individuo} & \textbf{CI} ($X$) & \textbf{Puntaje} ($Y$)\\
\hline
1 & 100  & 43 \\
2 & 100 & 42 \\
3 & 105 & 45 \\
4 & 103 & 44 \\
5 & 106 & 47 \\
6 & 115 & 45 \\
7 & 120 & 50 \\
\hline
\end{tabular}
\end{center}

el modelo de regresión planteado por los investigadores está dado por:
$$Puntaje = \beta_0 + \beta_1 CI + u$$

\begin{enumerate}
      \item (10 puntos) Determine los estimadores de Mínimos Cuadrados Ordinarios (MCO) 
$\hat{\beta}_0$ y $\hat{\beta}_1$ en el modelo que relaciona el 
\textit{Puntaje} en función del \textit{Coeficiente Intelectual (CI)}. 
Interprete el significado económico/estadístico de cada uno de estos estimadores.



{\color{blue} Calculamos $\Bar{X}$ y $\Bar{Y}$
$$\Bar{X} = (100+100+105+103+106+115+120)/7=107$$
$$\Bar{Y}=(43+42+45+44+47+45+50)/7=45,14$$
\begin{center}
\begin{tabular}{|c|c|c|c|c|c|c|}
\hline
\textbf{Individuo} & \textbf{CI} ($X$) & \textbf{Puntaje} ($Y$) & $X_i-\Bar{X}$ & $(Y_i - \Bar{Y})$ & $(X_i-\Bar{X})(Y_i-\Bar{Y})$ & $(X_i-\Bar{X})^2$\\
\hline
1 & 100  & 43 & -7 & -2,14  & 15    & 49\\
2 & 100 & 42 &  -7 & -3.14 & 22    &   49 \\
3 & 105 & 45 & -2  & -0.14   & 0.28  & 4  \\
4 & 103 & 44 & -4  & -1.14   & 4.56     & 16 \\
5 & 106 & 47 & -1  & 1.86   & -1.86    & 1  \\
6 & 115 & 45 &  8  & -0.14 & -1.12    & 64 \\
7 & 120 & 50 & 13  & 4.86    & 63.18    & 169 \\
\hline
total & & & & & 102 & 352 
\end{tabular}
\end{center}
por tanto $\beta_1$:
$$\hat{\beta_1} = \frac{102}{352}=0,29$$
y $\hat{\beta_0}$ sería:
$$\hat{\beta_0}= 45.14-0.29\cdot 107 = 14.11$$
\\ 
    El valor de $\hat{\beta_0}$ representa el intercepto de la RLS, mientras que el valor de $\hat{\beta_1}$ indica cuanto cambia Y, cuando X cambia en 1 unidad (pendiente).}

    \item (6 puntos) Calcule los valores predichos del \textit{Puntaje} a partir del modelo estimado.

\item (6 puntos) Represente gráficamente la \textbf{recta de regresión muestral} obtenida. 
En el mismo gráfico, marque los valores observados del \textit{Puntaje} 
y señale los \textbf{residuos} correspondientes a cada observación.

\begin{figure}[H]
    \centering
    \includegraphics[width=0.6\linewidth]{certamenes//Certamenes T2/graficoreg_CI.png}
\end{figure}

    \item (8 puntos) Calcule el coeficiente de determinación $R^2$ del modelo e interprete su valor.


 {\color{blue}
    Para calcular $R^2$, se debe calcular SST y SSR, usando la tabla tenemos:
\begin{center}
\begin{tabular}{|c|c|c|c|c|c|c|c|}
\hline
\textbf{Individuo} & \textbf{CI} ($X$) & \textbf{Puntaje} ($Y$)  & $(Y_i - \Bar{Y})$ &$(Y_i - \Bar{Y})^2$ & $\hat{Y_i}$ & $\hat{u}_i=Y_i-\hat{Y_i}$ & $\hat{u_i}^2$ \\
\hline
1 & 100  & 43 & -2,14 & 4.58 & 43.14 & -0.14  &  0.02\\
2 & 100 & 42  & -3.14 & 9.86 & 43.14 & -1.14 &  1.30\\
3 & 105 & 45  & -0.14 & 0.02 & 44.59 & 0.41  &  0.17\\
4 & 103 & 44  & -1.14 & 1.30 & 44.01 & 0.01  &  0.00\\
5 & 106 & 46  & 1.86  & 3.46 & 44.88 & 2.12  & 4.49 \\
6 & 115 & 45  & -0.14 & 0.02 & 47.49 & -2.49 & 6.2\\
7 & 120 & 50  & 4.86  & 23.62 & 48.94 & 1.06  & 1.12\\
\hline
total & & & & 42.86 &  &  & 13.3 
\end{tabular}
\end{center}   
por tanto,
$$R^2 = 1-\frac{SSR}{SST} = 1-\frac{13.3}{42.86}=0.69$$

El modelo explica 69\% de la varianza de los puntajes.}

\end{enumerate}
\end{enumerate}

\section{Formulario}
\begin{itemize}

\item El intervalo de confianza de $100(1-\alpha)\%$ para la media poblacional $\mu$ cuando $\sigma$ es \emph{\textbf{conocido}}:
$\color{red}{\bar{x} \pm z_{\frac{\alpha}{2}}\frac{\sigma}{\sqrt{n}}}$

\item El intervalo de confianza de $100(1-\alpha)\%$ para la diferencia de medias
$\mu_x-\mu_y$ es:
${\color{red}(\bar{X}-\bar{Y}) \pm z_{\frac{\alpha}{2}}\sqrt{\frac{\sigma_x^2}{n_x}+\frac{\sigma_y^2}{n_y}}}$

\item  Estimadores MCO.
{\color{red}
$$\hat{\beta}_0 = \bar{y} - \hat{\beta}_1 \bar{x} \quad \quad   \hat{\beta}_1 = \frac{\sum_{i=1}^{n}(x_i - \bar{x})(y_i-\bar{y})}{\sum_{i=1}^{n}(x_i - \bar{x})^2}$$}
    
\item 
	{\color{red}	\[ SST = \sum_{i=1}^{n} (y_i - \bar{y})^2 \quad \quad  SSE = \sum_{i=1}^{n} (\hat{y}_i - \bar{y})^2 \quad \quad SSR = \sum_{i=1}^{n} \hat{u}_i^2  \]}
		
		

\begin{figure}[H]
    \centering
    \includegraphics[width=0.75\linewidth]{Figuras/normal.png}
\end{figure}




\end{document}